% ============================================================
% Chapter 5: Development
%   Section: Deployment
% Included from chapters/05-development/development.tex
% Source: docs/decisions/deployment.md; deploy/gcp/;
% docker/docker-compose.prod.yml, docker-compose.gcp.yml.
% ============================================================

\section{Deployment}
\label{sec:dev-deployment}

\paragraph{The production compose file.} Production runs the same
services as development with a layer of its own: only the API and the
interface are published, and only on the loopback interface; every
secret is required; containers restart unless stopped and rotate their
logs; the language model is a container of its own; and memory caps
come from measurements. On 1~October a full analysis ran inside the
development stack (75~s) and inside the production one (184~s, the model
on a CPU). The production run found that the development limits for the
model --- two concurrent calls, a 30~s timeout --- turned 3 of 4 CPU
verdicts into timeouts; production sets one call at a time and 180~s.
The X-Fact pilot later found the same for inference on a CPU (four
claims at once pushed embedding requests past their 30~s timeout), and
production now also runs two claims at once with a 120~s inference
timeout, as the evaluation runs do.

\paragraph{Where it runs.} One Google Cloud Compute Engine virtual
machine, an \texttt{e2-standard-4} (4 vCPU, 16~GB) in Madrid, runs the
whole compose stack. It was sized from the production stack's
measurements (about 9.6~GiB resident) and from the model's time per
claim on a CPU. Cloud Run was ruled out: the embedded vector store takes
an exclusive file lock, jobs live in memory, work runs in background
threads, and about 6~GiB of models must stay loaded, so each service
would have needed a redesign.

\paragraph{Infrastructure as code.} Terraform creates the project's
APIs, a service account for the machine that can read only its four
secrets and write logs and metrics, the four secrets in Secret Manager,
a network of its own whose only open port is HTTPS, the machine itself
with a boot script, a daily disk snapshot kept seven days and, when an
address is given, an uptime check and a search-health alert by email.

\paragraph{The front door.} Caddy serves the site over HTTPS only, with
a password on every page generated by Terraform (only its hash reaches
the machine), because the interface's proxies attach the API key. The
reader view is the exception and is public. With a domain, Caddy
obtains a certificate; without one, only the health endpoint is public
and the site is reached through an authenticated tunnel.

\paragraph{Checked without an account.} The Terraform is applied to
Floci, a local Google Cloud emulator, by one command of the project's
check: 32 resources are created, the generated password opens the
rendered hash, and a second plan is empty. Floci runs no guest, so the
machine's boot has never been executed; a first real deployment needs
a billed project and is described step by step in the repository. Two
defects were found on the way: shell scripts checked out on Windows
with CRLF line endings, which would have broken the machine's boot
script, and a password hash whose dollar signs were expanded by the
shell.

\paragraph{Observability.} Logs are one JSON object per line; a
metrics endpoint in Prometheus' format counts stages, outcomes, verdicts
and empty or failed searches; and a health endpoint asks every
dependency. On Google Cloud an agent ships both to Cloud Logging and
managed Prometheus.
